
\medskip
\phantomsection\label{prop:sga3-VIB-1.4}%
\noindent\textbf{Proposition 1.4.---}
\emph{Let \(A\) be an Artinian local ring, and let \(u:G\to H\)
be a quasi-compact morphism between \(A\)-groups locally of finite
type. The following assertions are equivalent:
\begin{enumerate}
\item[\textup{\SGAInlineRefTarget{prop:sga3-VIB-1.4:item:i}(i)}] \(u\) is proper;
\item[\textup{\SGAInlineRefTarget{prop:sga3-VIB-1.4:item:ii}(ii)}] there exists \(h\in H\) such that the fiber
  \(u^{-1}(h)\) is nonempty and proper over \(\kappa(h)\);
\item[\textup{\SGAInlineRefTarget{prop:sga3-VIB-1.4:item:iii}(iii)}] \(\ker u\) is proper over \(A\).
\end{enumerate}}

\emph{Proof.}
It is clear that \textup{\SGARefLink{prop:sga3-VIB-1.4:item:i}{(i)}} implies \textup{\SGARefLink{prop:sga3-VIB-1.4:item:iii}{(iii)}}, and that
\textup{\SGARefLink{prop:sga3-VIB-1.4:item:iii}{(iii)}} implies \textup{\SGARefLink{prop:sga3-VIB-1.4:item:ii}{(ii)}}. The hypotheses imply that
\(u\) is of finite type; and since \(G\) is separated
(Expos\'e~VI\(_A\), 0.3), \(u\) is separated as well
(EGA~I, 5.5.1). It remains to prove that \textup{(ii)} implies
that \(u\) is universally closed. We may therefore suppose that
\(A\) equals its residue field \(k\).%
\footnote{\label{note:sga3-VIB-12}Editor's note: we have added the
following sentence.}
Let \(k'\) be an algebraic closure of \(\kappa(h)\), let
\[
  u':G'\longrightarrow H'
\]
be obtained from \(u\) by base change, and let \(h'\in H'\) be a
point above \(h\). Then
\[
  u'^{-1}(h')=u^{-1}(h)\times_{\kappa(h)}k'
\]
is nonempty and proper, and it is enough to prove that \(u'\) is
proper (EGA~IV\(_2\), 2.6.4). Thus we may suppose that \(k\) is
algebraically closed and \(h\in H(k)\).

By \hyperref[prop:sga3-VIB-1.2]{1.2}, \(u(G)\) is then the underlying set of a closed reduced
group subscheme of \(H\). Since every closed immersion is proper
(EGA~II, 5.4.2), we may suppose that \(u\) is surjective and that
\(H\) is reduced. The group \(G(k)\) then acts transitively on the
set of closed points of \(H\). Consequently, for every closed point
\(y\in H\), the fiber \(u^{-1}(y)\) is proper over \(\kappa(y)\).
By EGA~IV\(_3\), 9.6.1, the set of \(y\in H\) for which
\(u^{-1}(y)\) is not proper over \(\kappa(y)\) is locally
constructible. It contains no closed point, and is therefore empty
(cf.\ EGA~IV\(_3\), 10.3.1 and 10.4.7).

\footnote{\label{note:sga3-VIB-13}Editor's note: we have expanded
the original in what follows.}
Now consider the generic point \(\eta\) of \(H^0\). By what
precedes, the fiber
\[
  u^{-1}(\eta)=G\times_H\operatorname{Spec}\kappa(\eta)
\]
is proper over \(\kappa(\eta)\). Since \(H\) is reduced,
\(\kappa(\eta)=\mathcal O_{H,\eta}\). The latter ring is the direct
limit of the rings \(\mathcal O_H(V)\), as \(V\) ranges over the
nonempty open subsets of \(H^0\). It follows from
EGA~IV\(_3\), 8.1.2(a) and 8.10.5(xii), that there exists a
nonempty open subset \(V\subset H^0\) such that the restriction of
\(u\) over \(V\) is proper. The translates \(g\cdot V\), for
\(g\in G(k)\), form an open covering of \(H\); and for every
\(g\in G(k)\), the restriction of \(u\) over \(g\cdot V\) is
proper. Hence \(u\) is proper (cf.\ EGA~II, 5.4.1).

\medskip
\phantomsection\label{cor:sga3-VIB-1.4.1}%
\noindent\textbf{Corollary 1.4.1.---}
\emph{Let \(A\) be an Artinian local ring, and let \(u:G\to H\)
be a morphism between \(A\)-groups locally of finite type. The
following assertions are equivalent:
\begin{enumerate}
\item[\textup{(i)}] \(u\) is locally quasi-finite;
\item[\textup{(ii)}] \(u\) is quasi-finite at some point;
\item[\textup{(iii)}] \(\ker u\) is discrete;
\item[\textup{\SGAInlineRefTarget{cor:sga3-VIB-1.4.1:item:iv}(iv)}] the restriction of \(u\) to each connected
  component of \(G\) is finite.
\end{enumerate}
If \(u\) is quasi-compact, these assertions are also equivalent to
\begin{enumerate}
\item[\textup{(v)}] \(u\) is finite.
\end{enumerate}}

It is clear that \textup{(iv)} implies \textup{(iii)}, that
\textup{(iii)} implies \textup{(ii)} (EGA~I, 6.4.4), and that,
when \(u\) is quasi-compact, \textup{(iv)} and \textup{(v)} are
equivalent. We have already seen in \hyperref[prop:sga3-VIB-1.3]{1.3} that \textup{\SGARefLink{concl:sga3-VIB-1.3-i}{(i)}} and
\textup{\SGARefLink{prop:sga3-VIB-1.3:item:ii}{(ii)}} are equivalent.

It remains to prove that \textup{\SGARefLink{concl:sga3-VIB-1.3-i}{(i)}} implies \textup{\SGARefLink{cor:sga3-VIB-1.4.1:item:iv}{(iv)}}. Let
\(C\) be a connected component of \(G\). Since \(C\) is of finite
type over \(A\) (Expos\'e~VI\(_A\), \SGARefLink{cor:sga3-VIA-2.4.1}{2.4.1}) and \(G,H\) are
separated (Expos\'e~VI\(_A\), \SGARefLink{sga3:VIA:0.3}{0.3}),%
\begingroup
\renewcommand{\thefootnote}{\fnsymbol{footnote}}%
\footnote[1]{\label{note:sga3-VIB-1.4.1-separated-xref}Translator's
note: the French source cites Expos\'e~VI\(_A\), 0.2 here; the
separation statement is \SGARefLink{sga3:VIA:0.3}{0.3}.}%
\endgroup{}
EGA~I, 5.5.1 and 6.3.4 show
that the restriction \(u'\) of \(u\) to \(C\) is separated and of
finite type.%
\footnote{\label{note:sga3-VIB-14}Editor's note: we have expanded
the original in what follows.}
Its fibers are discrete, so \(u'\) is quasi-finite
(cf.\ EGA~II, 6.2.2). Since every quasi-finite proper morphism is
finite (cf.\ EGA~III\(_1\), 4.4.2), it is enough to prove that
\(u'\) is universally closed.

For this, we may suppose that \(A\) equals its residue field \(k\).
By \(\mathrm{fpqc}\) descent (cf.\ EGA~IV\(_2\), 2.6.4), it is
enough to prove that
\[
  u'\otimes_k\overline k
\]
is universally closed, where \(\overline k\) is an algebraic
closure of \(k\). Since \(C\) is of finite type over \(k\), the
scheme \(C\otimes_k\overline k\) is the disjoint union of finitely
many connected components \(C'_1,\ldots,C'_n\) of
\(G\otimes_k\overline k\), and it is enough to prove the assertion
for each \(C'_i\). We are thus reduced to the case in which \(k\)
is algebraically closed.

Let \(g\) be a closed point of \(C\), and let
\[
  u^0:G^0\longrightarrow H
\]
be the restriction of \(u\) to \(G^0\). Then
\[
  u'=r_{u(g)}\circ u^0\circ r_{g^{-1}},
\]
where \(r_g\) denotes right translation by \(g\). Thus, to prove
that \(u'\) is proper, it is enough to prove that \(u^0\) is
proper. By hypothesis \(u\) is locally quasi-finite, so the
nonempty fiber \(\ker u\) is discrete. We have seen that \(u^0\)
is of finite type, hence its fiber \(\ker u^0\) is finite
(cf.\ EGA~II, 6.2.2), and is therefore nonempty and proper.
\hyperref[prop:sga3-VIB-1.4]{Proposition~1.4} now shows that \(u^0\) is proper.

\medskip
\phantomsection\label{cor:sga3-VIB-1.4.2}%
\noindent\textbf{Corollary 1.4.2.---}
\emph{Let \(A\) be an Artinian local ring, and let \(u:G\to H\)
be a quasi-compact morphism between \(A\)-groups locally of finite
type. The following assertions are equivalent:%
\begin{enumerate}
\item[\textup{\SGAInlineRefTarget{cor:sga3-VIB-1.4.2:item:i}(i)}] \(u\) is a closed immersion;
\item[\textup{\SGAInlineRefTarget{cor:sga3-VIB-1.4.2:item:ii}(ii)}] \(u\) is a monomorphism;
\item[\textup{\SGAInlineRefTarget{cor:sga3-VIB-1.4.2:item:iii}(iii)}] \(\ker u\) is trivial, that is, isomorphic
  to the unit \(A\)-group.
\end{enumerate}
In particular, every group subscheme%
of \(H\) is closed.}

\begingroup
\renewcommand{\thefootnote}{\fnsymbol{footnote}}%
\footnote[1]{\label{note:sga3-VIB-1.4.2-unit-base}Translator's note:
the French source calls the trivial kernel the ``unit \(k\)-group''
both in the statement and in the proof. Here the base is the
Artinian local ring \(A\), and \(\ker u\) is an \(A\)-group; the
unit object required by the monomorphism criterion is therefore
\(\operatorname{Spec}A\), the unit \(A\)-group.}%
\endgroup{}
It is clear that \textup{\SGARefLink{cor:sga3-VIB-1.4.2:item:i}{(i)}} implies \textup{\SGARefLink{cor:sga3-VIB-1.4.2:item:ii}{(ii)}}. Considering
the functors represented by \(G\) and \(H\), one sees immediately
that \textup{\SGARefLink{cor:sga3-VIB-1.4.2:item:ii}{(ii)}} and \textup{\SGARefLink{cor:sga3-VIB-1.4.2:item:iii}{(iii)}} are equivalent. Finally, if
\(\ker u\) is the unit \(A\)-group, it is a proper nonempty fiber.
\hyperref[prop:sga3-VIB-1.4]{Proposition~1.4} therefore shows that \(u\) is a proper
monomorphism; it is of finite presentation because \(H\) is
locally Noetherian (EGA~IV\(_1\), 1.6.1), and hence is a closed
immersion (EGA~IV\(_3\), 8.11.5).

The final assertion follows because \(H\) is locally Noetherian,
so every immersion \(G\to H\) is quasi-compact (EGA~I, 6.6.4).
